\begin{figure}[t]
\begin{minted}[fontsize=\small,frame=single]{haskell}
-- Quipper teleportation example
teleport :: (Qubit, Qubit, Qubit) -> Circ (Qubit, Qubit, Qubit)
teleport (alice, bob, msg) = do
  -- Create entanglement between alice and bob
  alice <- hadamard alice
  bob <- qnot bob `controlled` alice
  
  -- Alice's Bell measurement
  msg <- qnot msg `controlled` alice
  msg <- hadamard msg
  
  -- Measure and get classical bits
  (m1, msg) <- measure msg
  (m2, alice) <- measure alice
  
  -- Classical control: conditionally apply corrections to bob
  bob <- controlled m2 (gate_X bob)
  bob <- controlled m1 (gate_Z bob)
  
  return (alice, bob, msg)
\end{minted}
\caption{Quantum teleportation in Quipper. The \texttt{controlled} keyword behaves differently for quantum vs. classical bits: when \texttt{m1} and \texttt{m2} are classical measurement results, gates are included in the circuit only if the bit equals 1. Dynamic lifting allows measurement results to guide circuit generation.}
\label{fig:quipper-teleport}
\end{figure}
